\section{Discussion}
\label{sec:discussion}

Read together, the comparisons say that a model's sensitivity to an earlier value tracks how that value is attached to an entity in the text, not what the value used to mean. A value that was the queried attribute and has since been updated, and a value that was some other attribute of the same entity and was never updated, differ by at most a third of a nat, although only the first is contradicted by the current state. A value placed in a sentence with the entity, even a note described as unrelated, is several nats more influential than either. A value attached to no one is credited to whichever entity shares its position. The word \lit{Previously} turns the influence down but not off. For the problem that motivated this work---whether models retain obsolete assignments---the answer from these scores is that the residual sensitivity we measure is not specific to obsolescence: the same sensitivity arises for values that were never assigned to the queried attribute at all.

The positional result connects our short prompts to findings in much longer ones. \citet{liu2024lost} and \citet{yang2026conflicts} report that where evidence appears changes how models use it, and \citet{gurarieh2026mixing} describe a positional mechanism for retrieving bound entities. Six-line contexts already show order deciding which entity an unattached value is credited to. A practical consequence for evaluation is that averaging over mention orders can produce null effects that conceal large opposing ones, as the near-zero overall relevance of unassigned values does here.

One gap remains open. Entity-mention relevance exceeds superseded relevance in Experiment 1, where only the superseded construction carries \lit{Previously}, and the gap persists in Experiment 2 when neither carries it: without the marker, $R$ is 7.304 [6.629, 7.966] for entity mention and 3.738 [3.536, 3.934] for superseded in Qwen, and 3.313 [2.975, 3.669] and 2.631 [2.409, 2.863] in Gemma. The two constructions still differ in predicate and word order. One property they differ in is the distance between the entity name and the value: in \lit{Nora mentioned amber}, the value follows the name after one word, while in \lit{Nora’s badge was amber} it follows the attribute and the copula. Given that relative position governs unattached values, we hypothesize that name--value proximity explains the construction gap. This predicts that moving the value closer to the name raises relevance within a construction, and that the gap shrinks at matched distance. [TODO: run P1.1 name--value distance test.]

The measure has a deliberately narrow scope. Scores are bounded masses of fixed answer prefixes, not probabilities of generated answers, and the models answer at ceiling: the current value is top-ranked in every superseded trial. At ceiling, the earlier value does not surface in generated answers: in a variant of Experiment 1 with two to six added distractor entities, neither model answered any of 1,152 superseded prompts with the earlier value (Appendix~\ref{app:harder}). Whether the score shifts change generated answers below ceiling is untested. Adding distractors did not lower accuracy enough to test it---Qwen stayed above 99.6\% at every level, which triggered that study's prespecified stopping rule. [TODO: run P1.3 harder-task redesign.] The experiments use two models, English prompts with two target entities, an eight-value vocabulary, and one set of sentence templates per experiment; generalization beyond these is untested.
